%% glossary entries
%\gls{trachytopes}
%\gls{maths} 

\newglossaryentry{trachytopes}
{
	name=trachytopes,
	description={Thrachytopes specify the bed roughness and flow resistance on a sub-grid level by defining and using various land use or roughness/resistance classes, further referred to as trachytopes after the Greek word $\tau \negthinspace \rho \alpha \chi \acute{\upsilon} \tau \eta \varsigma$ for roughness \citep{fm_manual}} 
}

\newglossaryentry{ploeg}
{
	name={ploeg},
    description={Mechanische ploeg van \si{6}{m} breed},
}

\newglossaryentry{ploegproductie}
{
	name={ploegproductie},
	description={De ploegproductie wordt in MvO (2025) gedefinieerd als \si{m^3} verplaatste bodemmateriaal ten opzichte van de geleverde inspanning. De geleverde inspanning is het aantal meters dat is gevaren met de ploeg. De ploegproductie vormt een kengetal afhankelijk van locatie}
}


\newglossaryentry{OLA}
{
	name={OLA},
	description={Overeengekomen Lage Afvoer. De OLA is de afvoer die
		gemiddeld 20 dagen per jaar wordt onderschreden. Voor de
		Rijntaken is deze afvoer momenteel \SI{1020}{m^3/s}.}
}

\newglossaryentry{OLR}
{
	name={OLR},
	description={Overeengekomen Lage Rivierstand. is een hoogte voor de ongestuwde delen van de Rijntakken, gebaseerd op de Overeengekomen Lage Afvoer (OLA). Het vormt een referentievlak ten opzichte waarvan diepten worden aangegeven  t.b.v. de scheepvaart.}
}

\newglossaryentry{BRV}
{
	name={BRV},
	description={Baggerreferentievlak. Een specifiek voor het baggeren gedefinieerd referentievlak dat hellend ligt ten opzichte van het N.A.P. en jaarlijks wordt vastgesteld. Het BRV is gerelateerd aan het diepte criterium: OLR - \SI{2.8}{m} (voor Boven-Rijn en Waal, Vanaf Waal rivierkilometer 918 neemt dit toe, zie bijlage 9 in RBK 6). Deze is in de Legger Rijkswaterstaatswerken vastgelegd.}
}

\newglossaryentry{RBK6}
{
	name={RBK 6},
	description={Rivierkundig beoordelingskader 6. \cite{RWS23}}
}


\newglossaryentry{NAP}
{
	name={N.A.P.},
	description={Normaal Amsterdams Peil}
}

\newglossaryentry{volume_above}
{
	name={volume above},
	description={Het volume dat relatief tot de nulmeting boven nul ligt. De nulmeting wordt ook wel de inpeiling genoemd.}
}

\newglossaryentry{volume_below}
{
	name={volume below},
	description={Het volume dat relatief tot de nulmeting onder nul ligt. Bij de bepaling van het onderhoudsvolume wordt de volumes below bij elke meting opgeteld tot één volume.}
}

\newglossaryentry{nulmeting}
{
	name={nulmeting},
	description={De nulmeting is de referentie waaraan de volumes below en above zijn bepaald. Het wordt de nulmeting genoemd omdat referentie verwarrend kan zijn met het BRV. De nulmeting is een vlak en is een direct resultaat van de inpeiling.}
}

\newglossaryentry{doelvolume}
{
	name={doelvolume},
	description={Het volume wat boven het BRV ligt wat voor de ondiepte zorgt. Het doelvolume kan worden benaderd door de inpeiling en het te behalen BRV van elkaar af te trekken.}
}

\newglossaryentry{onderhoudsinspanning}
{
	name={onderhoudsinspanning},
	description={De benodigde inspanning om de ondiepte te verwijderen. De inspanning uitgedrukt in volume hoeft niet gelijk te staan aan het doelvolume.}
}

\newglossaryentry{inspanningsvolume}
{
	name={inspanningsvolume},
	description={Het inspanningsvolume resulteerd uit de inspanning. De inspanning wordt geleverd waarbij er sediment wordt verplaatst. Het verplaatste sediment is de inspanningsvolume. Het inspanningsvolume is onafhankelijk van het doel volume. In veel gevallen is het inspanningsvolume veel groter dan het doel volume. Het inspanningsvolume kan nooit kleiner zijn dan het doel volume.}
}

\newglossaryentry{raai}
{
	name={raai},
	description={Een passage van een ploegschip in het te ploegen veld.}
}

\newglossaryentry{peilfrequentie}
{
	name={peilfrequentie},
	description={Het aantal gevaren raaien per tussen opeenvolgende multibeammetingen. Bijvoorbeeld: voor locatie 1 is er gestart met een peilfrequentie van één meting per raai. Daarna is overgegaan naar een meting na elke 4e raai (ofwel en peilfrequentie van 1 op 4). Deze term dient niet verward te worden met de \gls{bemonsteringsfrequentie}.}
}

\newglossaryentry{bemonsteringsfrequentie}
{
	name={bemonsteringsfrequentie},
	description={De frequentie van de van het meetinstrument. Bijvoorbeeld bij een multibeam echosounder \SI{200}{kHz} of \SI{400}{kHz}.}
}

\newglossaryentry{baggerproductie}
{
	name={baggerproductie},
	description={}
}

\newglossaryentry{ploegvak}
{
	name={ploegvak},
	description={Een oppervlak waarbinnen ploegen plaatsvindt.}
}
